\documentclass[11pt,aspectratio=169]{beamer}

\usetheme{Madrid}
\usecolortheme{default}

\usepackage{amsmath,amsfonts,amssymb}
\usepackage{bm}
\usepackage{graphicx}
\usepackage{algorithm}
\usepackage{algpseudocode}
\usepackage{xspace}

% Theorem environments
\newtheorem{claim}[theorem]{Claim}
\newtheorem{prop}[theorem]{Proposition}

% Custom commands based on the paper
\newcommand{\lis}{{\tt lis}}
\newcommand{\loss}{{\tt loss}}
\newcommand{\est}{{\tt est}}
\newcommand{\bx}{Box}
\newcommand{\errorcont}{\Psi}
\newcommand{\RR}{\mathbb{R}}
\newcommand{\bmu}{\overline{\mu}}
% \newcommand{\bx}{Box}
\newcommand{\green}{Green}
\newcommand{\red}{Red}
\newcommand{\classifyb}{ClassifyBox}
\newcommand{\approxLIS}{\approxlis}
\newcommand{\sample}{Sample}

\newcommand{\len}{len}
\newcommand{\wt}{wt}
\newcommand{\st}{st}
\newcommand{\w}[1]{w(#1)}
\newcommand{\m}[1]{m(#1)}
\newcommand{\find}{Find}
\newcommand{\safe}{Safe}
\newcommand{\unsplit}{Unsplittable}
\newcommand{\tree}{tree}
\newcommand{\imp}{Improved}
\newcommand{\sparse}{Sparse}
% \newcommand{\lis}{{\tt lis}}
% \newcommand{\est}{{\tt est}}
% \newcommand{\loss}{{\tt loss}}
\newcommand{\alis}{{\tt alis}}
\newcommand{\aloss}{{\tt aloss}}
\newcommand{\lisloss}{\lambda}
\newcommand{\lossin}{{\tt loss_in}}
\newcommand{\inside}{{\tt in}}
\newcommand{\outside}{{\tt out}}
\newcommand{\phase}{{\tt ph}}
\newcommand{\In}{In}
\newcommand{\Out}{Out}
\newcommand{\leftedge}{{\bf{\rm LeftEdge}}}
\newcommand{\approxloss}{{\rm{\bf ApproxLoss}}}
\newcommand{\terminalbox}{{\rm{\bf TerminalBox}}}
\newcommand{\grid}{Grid}
\newcommand{\good}{Good}
\newcommand{\bad}{Bad}
\newcommand{\Bad}{Bad}
\newcommand{\Good}{Good}
\newcommand{\size}{Size}
\newcommand{\dead}{{\tt disposable}}
\newcommand{\spar}{{\tt sparse}}

\newcommand{\deltapar}{\overline{\delta}}
\newcommand{\deltaval}{1/\errorcont}
\newcommand{\err}{\xi}
\newcommand{\gridprecval}{\frac{1}{(C_2\errorcont)^4}}
\newcommand{\kappaest}{\widehat{\kappa}}
\newcommand{\suff}{{21}}
\newcommand{\sszval}{100\errorcont^3}
\newcommand{\taint}{\eta}
\newcommand{\taintval}{1/10\errorcont}
\newcommand{\taintprob}{\phi}
\newcommand{\taupar}{\overline{\tau}}


\newcommand{\lin}[2]{Line(#1^{(#2)})}
\newcommand{\move}[3]{#1^{(#2:#3)}}
\newcommand{\leftpar}{leftpar}
\newcommand{\nonleftpar}{nonleftpar}
\newcommand{\nlpath}{path}
\newcommand{\bdim}{dim}
\newcommand{\cyl}{cyl}
\newcommand{\pr}{{\rm Pr}}
\newcommand{\Sift}{{\em Sift}}
\newcommand{\Rep}{Rep}
\newcommand{\Linesift}{{\em Linesift}}
\newcommand{\Build}{{\em Build}}
\newcommand{\poly}{\textrm{poly}}
\newcommand{\p}{\mathcal{P}}
\newcommand{\mymod}{\textrm{mod} \ }
\newcommand{\otilde}{\widetilde{O}}
%\newcommand{\qed}{\hfill $\Box$}
\newcommand{\eps}{\varepsilon}
\newcommand{\fu}{\varphi}
\newcommand{\restr}[1]{{|_{{\textstyle #1}}}}
\newcommand{\proj}{\mbox{\rm proj}}
\newcommand{\vol}{\mbox{\tt vol}\,}
\newcommand{\area}{\mbox{\tt area}\,}
\newcommand{\conv}{\mbox{\tt conv}\,}
\newcommand{\diam}{\hbox{\tt diam}\,}
\newcommand{\hdisc}{\mbox{\rm herdisc}}
\newcommand{\ldisc}{\mbox{\rm lindisc}}
\newcommand{\EX}{\hbox{\bf E}}
\newcommand{\prob}{{\rm Prob}}
\newcommand{\proofend}{{\medskip\medskip}}
%\newcommand{\proof}{{\noindent\bf Proof: }}
\newcommand{\reals}{{\rm I\!\hspace{-0.025em} R}}
\newcommand{\dist}{\hbox{dist}}
\newcommand{\defeq}{\mbox{\,$\stackrel{\rm def}{=}$\,}}
\newcommand{\boxalg}[1]
{\begin{center}\fbox{\parbox{\columnwidth}{\tt
\begin{tabbing}
\=mm\=mm\=mm\=mm\=mm\=mm\=mm\=mm\=mm\kill
#1
\end{tabbing} } } \end{center} }
%%%%Mike's commands August 2010
\newcommand{\mass}{{\bf {\rm mass}}}
\newcommand{\viol}{{\bf {\rm  viol}}}
\newcommand{\strip}[2]{#1|#2}
\newcommand{\width}{{\bf {\rm  w}}}
\newcommand{\longpath}{{\bf {\rm Long-Path}}}
\newcommand{\errorprob}{\eps}
\newcommand{\spliterror}{\sigma}
\newcommand{\griderror}{\xi}
\newcommand{\gridprec}{\alpha}
\newcommand{\thresh}{\theta}
% \newcommand{\errorcont}{\Psi}
\newcommand{\numclaims}{6}
\newcommand{\neterror}{\xi}
\newcommand{\netprec}{\alpha}
\newcommand{\adderror}{\alpha}
%%%lossA  is the error introduced by step A in the analysis
\newcommand{\lossone}{\lambda_1}  
\newcommand{\losstwo}{\lambda_2}
\newcommand{\lossthree}{\lambda_3}
\newcommand{\lossfour}{\lambda_4}
\newcommand{\lossfive}{\lambda_5}
\newcommand{\losssix}{\lambda_6}
\newcommand{\lossseven}{\lambda_7}
\newcommand{\approxzero}{A_0}
\newcommand{\approxone}{A_1}
\newcommand{\approxtwo}{A_2}
\newcommand{\approxthree}{A_3}
\newcommand{\approxfour}{A_4}
\newcommand{\approxfive}{A_5}
\newcommand{\approxsix}{A_6}
\newcommand{\approxseven}{A_7}
\newcommand{\errorterm}{\zeta}
\newcommand{\errorzero}{\varepsilon_0}
\newcommand{\errorone}{\varepsilon_1}
\newcommand{\errortwo}{\varepsilon_2}
\newcommand{\errorthree}{\varepsilon_3}
\newcommand{\errorfour}{\varepsilon_4}
\newcommand{\errorfive}{\varepsilon_5}
\newcommand{\errorsix}{\varepsilon_6}
\newcommand{\errorseven}{\varepsilon_7}
\newcommand{\errorA}{\varepsilon_A}
\newcommand{\algcomment}[1]{$\backslash * \backslash$ {#1}}

\newcommand{\ssize}{{\bf {\tt sample\_size}}}
\newcommand{\ssz}{\sigma}
\newcommand{\rvmean}{\lambda}
\newcommand{\samplemean}{{\bf {\rm sample-mean}}}
\newcommand{\Estimate}{{\bf {\rm Estimate}}}
\newcommand{\criticalbox}{{\bf {\rm Critical-Box}}}
\newcommand{\constructcall}{{\bf {\rm Construct-Call}}}
\newcommand{\irrelevant}{{\bf {\rm Irrelevant}}}
\newcommand{\rightedge}{{\bf {\rm rightedge}}}
\newcommand{\classify}{{\rm {\bf Classify}}}
\newcommand{\LIS}{{\rm {\bf LIS}}}
\newcommand{\basicmain}{{\rm {\bf BasicMain}}}
\newcommand{\improvedmain}{{\rm {\bf ImprovedMain}}}
\newcommand{\approxlis}{{\rm {\bf ApproxLIS}}}
\newcommand{\approxlisfin}{{\rm {\bf FinalApproxLIS}}}
\newcommand{\critbox}{{\rm {\bf CriticalBox}}}
\newcommand{\leafbox}{{\sc LeafBox}}
\newcommand{\nextbox}{{\sc NextBox}}
\newcommand{\splitbox}{{\bf{\rm SplitBox}}}
\newcommand{\childbox}{{\rm {\bf ChildBox}}}
\newcommand{\findsplitter}{{\rm {\bf FindSplitter}}}
\newcommand{\splitterfound}{splitter\_found}
\newcommand{\buildgrid}{{\rm {\bf BuildGrid}}}
\newcommand{\buildnet}{{\rm {\bf BuildNet}}}
\newcommand{\lischain}{{\rm {\bf LIS.Chain}}}
\newcommand{\gridpath}{{\rm {\bf GridPath}}}
\newcommand{\gridchain}{{\rm {\bf GridChain}}}
\newcommand{\point}[2]{\langle #1,#2 \rangle}
\newcommand{\expressionone}{E_1}
\newcommand{\expressiontwo}{E_2}
\newcommand{\expressionthree}{E_3}
\newcommand{\expressionfour}{E_4}
\newcommand{\expressionfive}{E_5}
\newcommand{\expressionsix}{E_6}
\newcommand{\expressionseven}{E_7}
\newcommand{\expressioneight}{E_8}
\newcommand{\maxval}{{\tt valbound}}
\newcommand{\liserror}{\nu}
\newcommand{\cP}{{\cal P}}
\newcommand{\cT}{{\cal T}}


\title[Estimating LIS in Polylog Time: Advanced Theory]{Estimating the Longest Increasing Sequence in Polylogarithmic Time}
\subtitle{Rigorous Analysis, Definitions, and Proofs}
\author{Based on the paper by M. Saks and C. Seshadhri}
\date{Advanced Seminar Presentation (2-3 Hours)}

\begin{document}

\begin{frame}
    \titlepage
\end{frame}

\begin{frame}{Outline}
    \tableofcontents[hideallsubsections]
\end{frame}

% =========================================================================
\section{Introduction and Formal Notation}
% =========================================================================

\begin{frame}{The Problem and Results}
    \begin{itemize}
        \item \textbf{Given:} An array represented as a function $f:[n] \rightarrow \mathbb{R}$.
        \item \textbf{Length of LIS:} $\lis_f$. \textbf{Distance to Monotonicity:} $\loss_f = n - \lis_f$.
        \item \textbf{Theorem 1.1:} There is a randomized algorithm computing $\est$ such that $|\est - \lis_f| \leq \delta n$ with probability $\geq 3/4$ in time $(1/\delta)^{O(1/\delta)} (\log n)^c$.
        \item \textbf{Theorem 1.2:} There is a $(1+\tau)$-approximation to the distance to monotonicity running in time $(1/(\epsilon_f\tau))^{O(1/\tau)}(\log n)^c$.
    \end{itemize}
    \vspace{0.5cm}
    \textit{This presentation focuses on the rigorous mathematical architecture proving these bounds.}
\end{frame}

\begin{frame}{Points and Partial Orders}
    \begin{definition}[Points and the Set $\mathcal{F}$]
        We map the array $f$ to points in $\mathbb{N}^2$. Let $F(x) = \langle x, f(x) \rangle$. We define $\mathcal{F} = \{F(x) : x \in (0, n]\}$.
    \end{definition}
    
    \begin{definition}[Relations between Points]
        For points $P, Q$:
        \begin{itemize}
            \item $P \preceq Q$: $x(P) \leq x(Q)$ and $y(P) \leq y(Q)$.
            \item $P \prec Q$: $x(P) < x(Q)$ and $y(P) \leq y(Q)$.
            \item $P \searrow Q$: $x(P) < x(Q)$ and $y(P) > y(Q)$ (a \textbf{violation}).
        \end{itemize}
        For $\mathcal{F}$-points, $P \sim Q$ denotes they are comparable (either $P \prec Q$ or $Q \prec P$). $P \not\sim Q$ denotes a violation.
    \end{definition}
\end{frame}

\begin{frame}{Regions Relative to a Point}
    For a given point $P$, we partition the plane $\mathbb{N}^2$ into four sets:
    \begin{itemize}
        \item $P^{SW} = \{Q \mid Q \preceq P\}$ (Southwest)
        \item $P^{NE} = \{Q \mid P \prec Q\}$ (Northeast)
        \item $P^{NW} = \{Q \mid x(Q) \leq x(P) \text{ and } y(Q) > y(P)\}$ (Northwest)
        \item $P^{SE} = \{Q \mid P \searrow Q\}$ (Southeast)
    \end{itemize}
    \vspace{0.5cm}
    \textbf{Note:} If $Q \in P^{NW} \cup P^{SE}$, then $Q$ and $P$ are in violation ($Q \not\sim P$).
\end{frame}

\begin{frame}{Boxes and Box Chains}
    \begin{definition}[Box]
        A box $\mathcal{B} = X(\mathcal{B}) \times Y(\mathcal{B})$ where $X(\mathcal{B})$ is an index interval $(x_L, x_R]$ and $Y(\mathcal{B})$ is a value interval $[y_B, y_T]$. The \textit{width} of $\mathcal{B}$ is $|X(\mathcal{B})| = x_R - x_L$.
    \end{definition}

    \begin{definition}[Box Chain]
        Boxes $\mathcal{B}_1 \triangleleft \mathcal{B}_2$ if the top-right of $\mathcal{B}_1$ is the bottom-left of $\mathcal{B}_2$. A sequence $\vec{\mathcal{T}} = (\mathcal{T}_1, \dots, \mathcal{T}_k)$ satisfying $\mathcal{T}_1 \triangleleft \cdots \triangleleft \mathcal{T}_k$ is a \textbf{box chain}.
    \end{definition}
    
    Any increasing sequence of points $P_0 \prec P_1 \prec \dots \prec P_k$ uniquely defines a box chain spanning from $P_0$ to $P_k$ via $\bx(P_{i-1}, P_i)$.
\end{frame}

\begin{frame}{Strips and Local LIS}
    \begin{definition}[Strips]
        A $\mathcal{B}$-strip is a subbox $\mathcal{S} \subseteq \mathcal{B}$ such that $Y(\mathcal{S}) = Y(\mathcal{B})$. It represents a vertical slice of $\mathcal{B}$. If $\mathcal{T} \subset \mathcal{B}$, $\mathcal{T}|\mathcal{B}$ denotes the $\mathcal{B}$-strip spanning the index interval of $\mathcal{T}$.
    \end{definition}
    
    \begin{definition}[Local \lis{} and \loss{}]
        For any box $\mathcal{B}$:
        \begin{itemize}
            \item $\lis(\mathcal{B})$ is the size of the longest $(F, \mathcal{B})$-increasing sequence.
            \item $\loss(\mathcal{B}) = |\mathcal{B} \cap \mathcal{F}| - \lis(\mathcal{B})$.
        \end{itemize}
    \end{definition}
\end{frame}

% =========================================================================
\section{Algorithmic Primitives: Splitters and Grids}
% =========================================================================

\begin{frame}{Splitters: Balance and Safeness}
    To simulate dividing and conquering the space, we search for indices to split a box $\mathcal{T} \subseteq \mathcal{B}$.
    
    \begin{definition}[$\rho$-Balanced]
        An index $s \in X(\mathcal{T})$ is $\rho$-balanced in $\mathcal{T}$ if $s - x_L(\mathcal{T}) \geq \rho \cdot \width(\mathcal{T})$ and $x_R(\mathcal{T}) - s \geq \rho \cdot \width(\mathcal{T})$.
    \end{definition}
    
    \begin{definition}[Safeness Parameters]
        Let $\mu \in (0,1)$, $L > 0$. Let $viol(s, \mathcal{S})$ be the number of points in $\mathcal{S} \cap \mathcal{F}$ violating $F(s)$. 
        \begin{itemize}
            \item We define $Z(s, \mathcal{S}) = viol(s, \mathcal{S}) - \mu |\mathcal{F} \cap \mathcal{S}|$.
            \item A strip $\mathcal{S}$ is $(\mu, L)$-accepting if $Z(s, \mathcal{S}) \leq L$.
        \end{itemize}
    \end{definition}
\end{frame}

\begin{frame}{Adequate Splitters}
    \begin{definition}[$(\mu, L, \rho)$-Adequate Splitter]
        Let $\mathcal{T} \subseteq \mathcal{B}$ be boxes. A splitter $s \in X(\mathcal{T})$ is $(\mu, L, \rho)$-adequate if:
        \begin{enumerate}
            \item $F(s) \in \mathcal{T}$.
            \item $s$ is $\rho$-balanced in $\mathcal{T}$.
            \item $s$ is $(\mu, L)$-safe for the strip $\mathcal{T}|\mathcal{B}$ (meaning every adjacent strip $\mathcal{S} \subseteq \mathcal{T}|\mathcal{B}$ is $(\mu, L)$-accepting).
        \end{enumerate}
    \end{definition}
    
    \begin{block}{$\findsplitter$ Procedure}
        Randomly samples $O(\log^2 n / \rho)$ balanced points. For each, it estimates $Z(s, \mathcal{S})$ via random sampling $approxZ$. By Hoeffding bounds, if an adequate splitter exists with frequency $\geq \rho$, it finds a valid splitter w.h.p.
    \end{block}
\end{frame}

\begin{frame}{The Dichotomy Lemma: Setup}
    \textbf{The most critical structural lemma of the paper.} If a box lacks good splitters, the LIS must lose a significant fraction of points.
    
    \begin{itemize}
        \item Let $\mathcal{R}$ be a box with a strip decomposition $\vec{\mathcal{S}}$. Let $s(x)$ be the width of the strip containing $x$.
        \item An index $x$ is \emph{safe} if it is a $(\mu, s(x))$-safe splitter for $\mathcal{R}$. Otherwise, \emph{unsafe}.
        \item Let $\vec{\mathcal{D}}$ be an arbitrary box-chain compatible with $\vec{\mathcal{S}}$.
        \item Let $S$ be the safe indices in $\vec{\mathcal{D}}^\cup$, and $U$ be the unsafe indices in $\vec{\mathcal{D}}^\cup$.
        \item Let $\chi = |\mathcal{R} \cap \mathcal{F}|$. $\chi^{in}$ points inside $\vec{\mathcal{D}}^\cup$, $\chi^{out}$ points outside.
    \end{itemize}
\end{frame}

\begin{frame}{The Dichotomy Lemma (Lemma 4.6)}
    \begin{block}{Lemma 4.6}
        Under the defined setup, the number of points outside the chain satisfies:
        $$\chi^{out} \geq \frac{\mu}{1-\mu} |U|$$
        Consequently, the number of points inside satisfies:
        $$\chi^{in} \leq (1-\mu)\chi + \mu |S|$$
    \end{block}
    \vspace{0.3cm}
    \textit{Significance:} If there are few safe points ($|S|$ is small), then $\chi^{in}$ is bounded strictly away from $\chi$. Any increasing sequence MUST miss a fraction of points proportional to $\mu$.
\end{frame}

\begin{frame}{Proof of the Dichotomy Lemma (1/3)}
    \begin{proof}
        For each unsafe index $x \in U$, there exists an adjacent strip $\mathcal{W}(x)$ that is $(\mu, s(x))$-rejecting. This means $viol(x, \mathcal{W}(x)) \geq \mu |\mathcal{F} \cap \mathcal{W}(x)| + s(x)$.
        
        Let $\mathcal{V}(x)$ be the set of violations with $F(x)$ contained in $\mathcal{W}(x)$ but \textit{not} in $\vec{\mathcal{D}}^\cup$.
        
        \textbf{Claim:} $|\mathcal{V}(x)| \geq \frac{\mu}{1-\mu} |\vec{\mathcal{D}}^\cup \cap \mathcal{W}(x) \cap \mathcal{F}|$.
        
        \textit{Proof of Claim:} Any violation with $x$ contained \textit{inside} $\vec{\mathcal{D}}^\cup$ must reside in the single box $\vec{\mathcal{D}}[x]$. The number of such points is at most $\width(\vec{\mathcal{D}}[x]) \leq s(x)$.
    \end{proof}
\end{frame}

\begin{frame}{Proof of the Dichotomy Lemma (2/3)}
    \begin{proof}
        Therefore, subtracting the internal violations, the external violations satisfy:
        $$|\mathcal{V}(x)| \geq \mu |\mathcal{F} \cap \mathcal{W}(x)| + s(x) - s(x) = \mu |\mathcal{F} \cap \mathcal{W}(x)|$$
        Since $\mathcal{V}(x)$ is disjoint from $\vec{\mathcal{D}}^\cup$, we have:
        $$|\vec{\mathcal{D}}^\cup \cap \mathcal{W}(x) \cap \mathcal{F}| \leq (1-\mu) |\mathcal{F} \cap \mathcal{W}(x)|$$
        Rearranging gives $|\mathcal{V}(x)| \geq \frac{\mu}{1-\mu} |\vec{\mathcal{D}}^\cup \cap \mathcal{W}(x) \cap \mathcal{F}|$. \qed (Claim)
    \end{proof}
\end{frame}

\begin{frame}{Proof of the Dichotomy Lemma (3/3)}
    \begin{proof}
        We greedily select a subset $L' \subseteq L$ (unsafe indices facing left) and $R' \subseteq R$ (unsafe facing right) such that the strips $\mathcal{W}(z)$ are pairwise disjoint, yet they cover all $F(x)$ for $x \in U$.
        
        The external violations $\mathcal{V}(x)$ lie in the Northwest or Southeast regions, completely disjoint from the box chain $\vec{\mathcal{D}}$. 
        $$\chi^{out} \geq \Big| \bigcup_{x \in L' \cup R'} \mathcal{V}(x) \Big| = \sum_{x \in L' \cup R'} |\mathcal{V}(x)|$$
        $$\geq \sum_{x \in L' \cup R'} \frac{\mu}{1-\mu} |\vec{\mathcal{D}}^\cup \cap \mathcal{W}(x) \cap \mathcal{F}| \geq \frac{\mu}{1-\mu} |U|$$
        Since $\chi^{in} = |U| + |S|$, substituting $\chi^{in} = \chi - \chi^{out}$ yields $\chi^{in} \leq (1-\mu)\chi + \mu|S|$.
    \end{proof}
\end{frame}

\begin{frame}{Value Nets and Grids}
    To compute boosting paths, we partition terminal boxes.
    
    \begin{definition}[Value Net]
        An $\alpha$-value net for $\mathcal{B}$ is a subset $V \subset Y(\mathcal{B})$ such that any subinterval $J$ containing at least $\alpha \cdot \width(\mathcal{B})$ points contains at least one element of $V$.
    \end{definition}
    
    \begin{definition}[$\alpha$-fine Grid]
        A grid $\Gamma$ is $\alpha$-fine for $\mathcal{B}$ if:
        \begin{itemize}
            \item The columns partition $X(\mathcal{B})$ into strips of width $\leq \alpha \cdot \width(\mathcal{B})$.
            \item The rows form an $(\alpha / |X(\Gamma)|)$-value net.
        \end{itemize}
    \end{definition}
    Constructed via random sampling in procedure $\buildgrid$.
\end{frame}

\begin{frame}{Grid Approximation Lemma}
    \begin{block}{Lemma 4.9}
        If $\Gamma$ is an $\alpha$-fine grid for $\mathcal{B}$, and $\mathcal{L}$ is an increasing sequence in $\mathcal{B}$, there exists a $D(\Gamma)$-chain $\vec{\mathcal{D}}$ such that:
        $$|\mathcal{L} \setminus \vec{\mathcal{D}}^\cup| \leq \alpha \cdot \width(\mathcal{B})$$
    \end{block}
    
    \begin{proof}[Proof Sketch]
        Let $x_1 < \dots < x_r$ be the columns of $\Gamma$. Let $Q_i$ be the largest point of $\mathcal{L}$ with index $\leq x_i$. Choose $P_i \in \Gamma$ in column $x_i$ as the lowest point $\succeq Q_i$. 
        
        Points of $\mathcal{L}$ lost outside $\bx(P_{i-1}, P_i)$ must lie in the $P_i^{SE}$ region bounded vertically between $Q_i$ and $P_i$. By the value net property, this gap contains at most $\alpha \cdot \width(\mathcal{B}) / r$ points. Summing over $r$ columns yields $\alpha \cdot \width(\mathcal{B})$.
    \end{proof}
\end{frame}

% =========================================================================
\section{The Basic Approximation Algorithm}
% =========================================================================

\begin{frame}{Structure of the Algorithm}
    The algorithm is deeply recursive. We evaluate $\approxlis_t(\mathcal{B})$ to approximate $\lis(\mathcal{B})$ at recursion depth $t$.
    
    \textbf{Core procedures:}
    \begin{enumerate}
        \item $\approxlis_t$: Randomly samples indices to estimate the size of the set $Good_t(\mathcal{B})$.
        \item $\classify_t(x, \mathcal{B})$: Determines if $x$ is $Good$ or $Bad$. Finds a $\critbox_t(x)$, then calls $\classify_{t-1}(x, \mathcal{C})$.
        \item $\critbox_t$: Uses $\terminalbox_t$ and $\gridchain_t$.
        \item $\terminalbox_t$: Iteratively splits $\mathcal{B}$ using $\findsplitter$ until failure, yielding a Terminal Chain.
        \item $\gridchain_t$: Builds an $\alpha$-fine grid in a terminal box, runs $\approxlis_{t-1}$ on grid boxes, and extracts the longest path.
    \end{enumerate}
\end{frame}

\begin{frame}{Defining the Approximation Goal}
    \begin{definition}[$(\tau, \delta)$-approximation]
        We say $\alis_t(\mathcal{B})$ is a $(\tau_t, \delta_t)$-approximation if:
        $$|\alis_t(\mathcal{B}) - \lis(\mathcal{B})| \leq \tau_t \loss(\mathcal{B}) + \delta_t \width(\mathcal{B})$$
    \end{definition}
    
    \begin{itemize}
        \item \textbf{Primary Error:} $\tau_t \loss(\mathcal{B})$. $\tau_t$ must shrink as $t$ increases.
        \item \textbf{Secondary Error:} $\delta_t \width(\mathcal{B})$. Accounts for sampling failure, grid alignment, and splitter violations.
    \end{itemize}
    \textbf{Target:} $\tau_t = 4/t$ and $\delta_t = t/\log n$.
\end{frame}

% =========================================================================
\section{Error Analysis: Primary and Secondary Error}
% =========================================================================

\begin{frame}{Transforming the Goal}
    Let $\nu_t(\mathcal{B}) = \lis(\mathcal{B}) - \alis_t(\mathcal{B})$. Because $\lis(\mathcal{B}) \geq |Good_t(\mathcal{B})|$, we define a discrepancy term $\Delta_t$:
    $$\Delta_t = \lis(\mathcal{B}) - |Good_t(\mathcal{B})| - \tau_t \loss(\mathcal{B})$$
    Our goal reduces to showing $\Delta_t \leq \delta_t \width(\mathcal{B}) - |\epsilon_1|$ (where $\epsilon_1$ is sampling error).
    
    We partition $\Delta_t$ over all terminal boxes $\mathcal{T}$ in the terminal chain $\vec{\mathcal{T}}$:
    $$\Delta_t(\mathcal{T}) = (1+\tau_t)|\mathcal{L}(\mathcal{T})| - |Good_t(\mathcal{B}) \cap \mathcal{T}| - \tau_t \beta(\mathcal{T})$$
    where $\mathcal{L}$ is the optimal LIS, and $\beta(\mathcal{T})$ is the number of $\mathcal{F}$-points in the vertical strip of $\mathcal{T}$.
\end{frame}

\begin{frame}{The Five Secondary Error Terms}
    To manage the inequality algebraically, we isolate five error components for each terminal box:
    \begin{itemize}
        \item $\epsilon_1 = \alis_t - |Good_t|$ (Top-level sampling error)
        \item $\epsilon_2(\mathcal{T}) = \sum_{\mathcal{C}} \alis_{t-1}(\mathcal{C}) - |Good_{t-1}(\mathcal{C})|$ (Recursive sampling error)
        \item $\epsilon_3(\mathcal{T}) = |\mathcal{L}^{in}(\mathcal{T})| - \sum_{\mathcal{E}} |\mathcal{E} \cap \mathcal{L}|$ (Grid alignment loss)
        \item $\epsilon_4(\mathcal{T}) = |\mathcal{L}^{out}(\mathcal{T})| - \mu_t \cdot \outside(\vec{\mathcal{E}}(\mathcal{T}))$ (Splitter violation loss)
        \item $\epsilon_5(\mathcal{T}) = \mu_t \cdot \aloss_{t-1}(\vec{\mathcal{E}}) - (1-\mu_t)\cdot \outside(\vec{\mathcal{E}})$ (Dichotomy lemma bound error)
    \end{itemize}
\end{frame}

\begin{frame}{Bounding the Primary Error (Claim 7.3)}
    By substituting the error terms into $\Delta_t(\mathcal{T})$, we obtain:
    \begin{align*}
        \Delta_t(\mathcal{T}) \leq &(1+\tau_t)|\mathcal{L}^{out}| - \tau_t \cdot \text{outside} - (1+\tau_t)\text{loss}(\vec{\mathcal{E}}) \\
        &+ \aloss_{t-1}(\vec{\mathcal{E}}) + \epsilon_2 + (1+\tau_t)\epsilon_3
    \end{align*}
    By the inductive hypothesis:
    $$\text{loss}(\vec{\mathcal{E}}) \geq \frac{\aloss_{t-1}}{1+\tau_{t-1}} - \frac{\delta_{t-1}\width}{1+\tau_{t-1}}$$
    Let $K_t = 1/(1+\tau_{t-1})$. Substituting this bounds the $\aloss$ term. We then use $\epsilon_4, \epsilon_5$ to eliminate $|\mathcal{L}^{out}|$ and $\aloss_{t-1}$.
\end{frame}

\begin{frame}{Setting the $\tau$ and $\mu$ parameters}
    The algebraic manipulation yields a coefficient for the $\text{outside}(\vec{\mathcal{E}}(\mathcal{T}))$ term:
    $$(1+\tau_t)\mu_t - \tau_t + [1 - K_t(1+\tau_t)](\mu_t^{-1} - 1)$$
    We MUST make this coefficient $\leq 0$. 
    
    \begin{itemize}
        \item We choose $\mu_t = (1-K_t)/2$ to minimize the quadratic form.
        \item This demands the condition:
        $$1 + \tau_t \geq \frac{4}{(1+K_t)^2} \implies \tau_t \geq \left( \frac{2\tau_{t-1} + 2}{\tau_{t-1} + 2} \right)^2 - 1$$
        \item Guessing $\tau_t = c/t$, we verify that $\tau_t = 4/t$ satisfies the recurrence! 
        \item Consequently, $\mu_t = 2/(t+3)$.
    \end{itemize}
\end{frame}

\begin{frame}{Bounding Secondary Error: $\epsilon_4$ (Lemma 7.7)}
    \begin{block}{Lemma 7.7}
        $$\sum_{\mathcal{T}} \epsilon_4(\mathcal{T}) \leq 2\alpha \width(\mathcal{B}) + 4 \sum_{P \in \cP^\circ} \gamma(P) \width(\mathcal{H}(P))$$
    \end{block}
    \begin{proof}[Proof Idea]
        $\epsilon_4$ bounds points of the LIS falling outside the terminal box but inside its vertical strip. These points form sets $\mathcal{L}^-$ and $\mathcal{L}^+$ bounding the top and bottom.
        
        Because the splitters defining the terminal box were $(\mu, 2\gamma)$-safe, the number of LIS points violating the corner splitters is strictly bounded by $\gamma \cdot \width(\mathcal{H}(P))$. Summing over the depth of the splitter tree yields the $O(\gamma \log n / \rho) \width(\mathcal{B})$ bound.
    \end{proof}
\end{frame}

\begin{frame}{Bounding Secondary Error: $\epsilon_5$ (Claim 7.8)}
    \begin{block}{Claim 7.8}
        If $\mathcal{T}$ is wide, $\epsilon_5(\mathcal{T}) \leq 4\rho \cdot \width(\mathcal{T})$.
    \end{block}
    \begin{proof}
        We apply the Dichotomy Lemma to the strip $\mathcal{T}|\mathcal{B}$ with grid decomposition $\Gamma$.
        $$(1-\mu_t) \cdot \outside(\vec{\mathcal{D}}) \geq \mu_t |U|$$
        Thus, $\epsilon_5(\mathcal{T}) \leq \mu_t (\aloss_{t-1} - |U|) \leq \mu_t |S|$.
        
        Since $\findsplitter$ FAILED on $\mathcal{T}$, the number of safe points $|S|$ is bounded by the number of unbalanced points $+1$, which is $\leq 3\rho \width(\mathcal{T}) + 1 \leq 4\rho \width(\mathcal{T})$.
    \end{proof}
\end{frame}

% =========================================================================
\section{The Improved Algorithm: Breaking the Logarithmic Barrier}
% =========================================================================

\begin{frame}{The Bottleneck in the Basic Algorithm}
    \begin{itemize}
        \item The running time is $O(\log n)^{O(1/\tau)}$.
        \item To drop the exponent of $\log n$ to a constant, we must reduce the recursion branching factor.
        \item The branching factor depends on the grid size $(1/\alpha)^{O(1)}$.
        \item In the basic algorithm, we needed $\alpha \leq \gamma \approx 1/\log n$. Why? Because applying the dichotomy lemma to bound $\epsilon_5$ required grid boxes to be thinner than the splitter allowance $\gamma$.
        \item If we increase $\gamma$, the splitter tree error $\epsilon_4 = \gamma \log n / \rho$ blows up!
    \end{itemize}
\end{frame}

\begin{frame}{The Innovation: Dynamic Phases in $\terminalbox$}
    Instead of terminating immediately when $\findsplitter$ fails, we introduce \textbf{Phases}.
    
    \begin{itemize}
        \item \textbf{Phase $j$}: We allow violation fraction $\gamma_j$ and balance $\rho_j$.
        \item If $\findsplitter$ fails, we DO NOT terminate.
        \item We increment $j$, exponentially increase $\gamma_j \to 16 \gamma_{j-1}$, and $\rho_j = \gamma_j^{1/4}$.
        \item We also raise the terminal width threshold $\theta$.
        \item We terminate only when $\width(\mathcal{T}) \leq \theta$.
    \end{itemize}
    \vspace{0.5cm}
    \textit{Consequence:} We decouple $\alpha$ from $\log n$. We can set $\alpha = \text{poly}(1/\errorcont)$, independent of $n$.
\end{frame}

\begin{frame}{The Improved $\terminalbox_t$}
    \begin{algorithmic}[1]
        \State $j \gets 0, \theta \gets \omega, \gamma \gets \gamma_0, \rho \gets \rho_0$
        \While{$\width(\mathcal{T}) > \theta$}
            \State Run $\findsplitter(\mathcal{T}, \dots, \gamma, \rho)$
            \If{splitter found}
                \State Split $\mathcal{T}$ and update boundaries.
            \Else
                \State \textit{// Phase fails. Relax parameters.}
                \State $\theta \gets \max(\theta, \gamma \width(\mathcal{T}) / \alpha)$
                \State $j \gets j+1$, $\gamma \gets \gamma_j$, $\rho \gets \rho_j$
            \EndIf
        \EndWhile
    \end{algorithmic}
\end{frame}

\begin{frame}{Revisiting Tainted Boxes (Lemma 9.7)}
    Because we shrunk the sample size $\sigma$ to be independent of $n$, we can no longer use a union bound over ALL boxes.
    
    \begin{definition}[Tainted Box-Level Pair]
        $(\mathcal{B}, t)$ is tainted if:
        1. Top-level sample fails: $|\alis_t - |Good_t|| > \Psi \width(\mathcal{B})$.
        OR 2. Recursive failure: A grid chain exists where the width of tainted children $> \Psi \width(\mathcal{B})$.
    \end{definition}
    
    \begin{lemma}[Lemma 9.7]
        The probability that $(\mathcal{B}, t)$ is tainted is $\leq \phi = \alpha^5 / C_2$.
    \end{lemma}
    \textit{Proof Sketch:} Proved via induction on $t$. Since $\alpha$-fine grids are disjoint, child failures are independent. A Chernoff bound on independent Bernoulli variables bounds the cascaded failure probability strictly below $\phi$.
\end{frame}

\begin{frame}{Fixing the $\epsilon_5$ Bound with Phases}
    In the basic algorithm, we bounded $|S|$ because $\findsplitter$ failed on $\mathcal{T}$. In the improved algorithm, $\cT$ terminated because its width dropped below $\theta$, not necessarily because a split failed.
    
    \begin{prop}[Proposition 9.9]
        For any wide terminal box $\cT$, there is a \textbf{multiphase box} $\mathcal{H}$ (an ancestor in the splitter tree) that failed $\findsplitter$ at phase $j-1$, such that every safe point in $\cT$ is $(\mu_t, \gamma_{j-1})$-safe for $\mathcal{H}$.
    \end{prop}
    
    This allows us to aggregate the $\epsilon_5$ error over the multi-phase boxes $\mathcal{H}$, bounding the total safe points by a geometric series $\sum_j 3 \rho_{j-1} \width(\mathcal{B})$, which converges to $O(\alpha^{1/4} \width(\mathcal{B}))$, completing the final error bound!
\end{frame}

\begin{frame}{Final Time Complexity}
    \begin{itemize}
        \item $\errorcont = \max(C_2, t_{\max}/\delta)$.
        \item Grid size is $O(1/\alpha) = O(\errorcont^4)$.
        \item Branching factor per recursive call is purely polynomial in $\errorcont$.
        \item Recursion depth $t_{\max} = O(1/\tau)$.
        \item \textbf{Total Time:} $(1/\errorcont)^{O(1/\tau)} \log^c n = (1/\delta\tau)^{O(1/\tau)} \log^c n$.
    \end{itemize}
    \vspace{0.5cm}
    \textbf{Conclusion:} A fully sublinear, additive $\delta n$ approximation to the LIS length!
\end{frame}

\begin{frame}
    \begin{center}
        \Huge Questions?
    \end{center}
    \vspace{1cm}
    \begin{center}
        \Large Thank you for your attention.
    \end{center}
\end{frame}

\end{document}
